%========================================
\chapter{网络层}
%========================================
\section{网络层服务}
网络层的主要功能是实现不同网络之间的互联，为分组选择合适的路由。

\subsection{两种服务模式}
\begin{table}[htbp]
\centering\small
\begin{tabular}{lcc}
\toprule
\textbf{特性} & \textbf{虚电路服务} & \textbf{数据报服务} \\
\midrule
连接建立 & 需要 & 不需要 \\
可靠性保证 & 网络保证 & 主机保证 \\
分组路径 & 固定 & 可变 \\
分组顺序 & 有序 & 可能乱序 \\
典型应用 & ATM & \textbf{因特网} \\
\bottomrule
\end{tabular}
\end{table}

\section{IPv4地址}
\subsection{地址格式}
32位二进制，用点分十进制表示：
\begin{center}
\texttt{10101100.00010000.00000001.00000001} $\Leftrightarrow$ \texttt{172.16.1.1}
\end{center}

\subsection{分类编址}
\begin{table}[htbp]
\centering\small
\begin{tabular}{ccccc}
\toprule
\textbf{类别} & \textbf{首位} & \textbf{地址范围} & \textbf{网络数} & \textbf{每网主机数} \\
\midrule
A & 0 & 1.0.0.0-126.x.x.x & 126 & $2^{24}-2$ \\
B & 10 & 128.0.0.0-191.255.x.x & $2^{14}$ & $2^{16}-2$ \\
C & 110 & 192.0.0.0-223.255.255.x & $2^{21}$ & 254 \\
D & 1110 & 224-239 & \multicolumn{2}{c}{多播地址} \\
E & 1111 & 240-255 & \multicolumn{2}{c}{保留地址} \\
\bottomrule
\end{tabular}
\end{table}

\subsection{特殊地址}
\begin{itemize}[leftmargin=2em,itemsep=0pt]
    \item \texttt{0.0.0.0}：本网络本主机
    \item \texttt{127.x.x.x}：环回测试地址
    \item \texttt{255.255.255.255}：本网段广播
    \item 主机号全0：网络地址
    \item 主机号全1：直接广播地址
\end{itemize}

\subsection{私有地址}
\begin{itemize}[leftmargin=2em,itemsep=0pt]
    \item A类：\texttt{10.0.0.0/8}
    \item B类：\texttt{172.16.0.0/12}
    \item C类：\texttt{192.168.0.0/16}
\end{itemize}

\section{子网划分}
\subsection{子网掩码}
子网掩码用于区分IP地址中的网络部分和主机部分。

\textbf{子网掩码的作用}：
\begin{itemize}[leftmargin=2em,itemsep=0pt]
    \item IP地址 AND 子网掩码 = 网络地址
    \item 用于判断两个IP地址是否在同一子网
\end{itemize}

\textbf{默认子网掩码}：
\begin{itemize}[leftmargin=2em,itemsep=0pt]
    \item A类：\texttt{255.0.0.0}（/8）
    \item B类：\texttt{255.255.0.0}（/16）
    \item C类：\texttt{255.255.255.0}（/24）
\end{itemize}

\begin{examplebox}[例题4.1 子网划分]
\textbf{题目}：C类地址\texttt{192.168.1.0}划分为6个子网，求子网掩码和每子网主机数。
\tcblower
需要$\lceil\log_2 6\rceil = 3$位子网号

子网掩码：\texttt{255.255.255.224}（/27）

主机位：$8-3=5$位，每子网主机数：$2^5-2=$ \textbf{30台}
\end{examplebox}

\section{CIDR无分类编址}
CIDR消除了传统的A、B、C类地址的界限。

\subsection{斜线记法}
\texttt{IP地址/前缀长度}，如\texttt{192.168.1.0/24}

\textbf{地址块大小}：$2^{32-n}$个地址（n为前缀长度）

\subsection{最长前缀匹配}
路由器转发时，选择与目的地址匹配的\textbf{最长前缀}的路由。

\begin{examplebox}[例题4.2 最长前缀匹配]
\textbf{题目}：目的地址\texttt{172.32.47.165}，路由表如下，转发给哪个路由器？
\begin{center}\small
\begin{tabular}{|c|c|}
\hline
目的网络 & 下一跳 \\
\hline
172.32.0.0/16 & R1 \\
172.32.32.0/19 & R2 \\
172.32.47.0/24 & R3 \\
0.0.0.0/0 & R4 \\
\hline
\end{tabular}
\end{center}
\tcblower
匹配分析：/16、/19、/24、/0都匹配

最长前缀是/24 $\Rightarrow$ 转发给\textbf{R3}
\end{examplebox}

\subsection{默认路由配置}
在路由器中配置默认路由通常使用以下命令：
\begin{center}
\texttt{ip route 0.0.0.0 0.0.0.0 <下一跳地址或出接口>}
\end{center}
\textbf{作用}：当目的地址在路由表中找不到匹配项时，按照默认路由转发。

\section{VLSM变长子网掩码}
VLSM允许在同一网络中使用不同长度的子网掩码。

\begin{examplebox}[例题4.3 VLSM地址分配]
\textbf{题目}：地址块\texttt{192.168.10.0/24}，分配给：销售部50台、技术部25台、财务部10台。
\tcblower
按需求从大到小分配：
\begin{enumerate}[itemsep=0pt]
    \item 销售部(50台)：用/26，分配\texttt{192.168.10.0/26}（62台）
    \item 技术部(25台)：用/27，分配\texttt{192.168.10.64/27}（30台）
    \item 财务部(10台)：用/28，分配\texttt{192.168.10.96/28}（14台）
\end{enumerate}
\end{examplebox}

\begin{examplebox}[例题4.4 路由聚合]
\textbf{题目}：将\texttt{10.1.192.0/24}、\texttt{10.1.193.0/24}、\texttt{10.1.194.0/24}、\texttt{10.1.195.0/24}聚合。
\tcblower
第三字节：192=11000000, 193=11000001, 194=11000010, 195=11000011

共同前缀：前6位相同(110000)

\textbf{聚合结果}：\texttt{10.1.192.0/22}
\end{examplebox}

\section{IP数据报格式}
\subsection{首部格式}
\begin{itemize}[leftmargin=2em,itemsep=0pt]
    \item \textbf{版本}（4位）：IPv4为4
    \item \textbf{首部长度}（4位）：以4字节为单位，最小5（20字节）
    \item \textbf{总长度}（16位）：整个数据报的长度，最大65535字节
    \item \textbf{标识、标志、片偏移}：用于分片和重组
    \item \textbf{生存时间TTL}（8位）：每经过一个路由器减1
    \item \textbf{协议}（8位）：TCP(6)、UDP(17)、ICMP(1)
    \item \textbf{首部校验和}（16位）：只校验首部
\end{itemize}

\subsection{IP分片}
当IP数据报长度超过链路的MTU时，需要进行分片。

\begin{examplebox}[例题4.5 IP分片]
\textbf{题目}：一个4000字节的IP数据报（首部20字节），经过MTU=1500字节的链路，如何分片？
\tcblower
数据部分：$4000 - 20 = 3980$字节

每个分片最大数据：$1500 - 20 = 1480$字节

分片情况：
\begin{center}\small
\begin{tabular}{|c|c|c|c|c|}
\hline
分片 & 数据长度 & 总长度 & 片偏移 & MF \\
\hline
1 & 1480 & 1500 & 0 & 1 \\
2 & 1480 & 1500 & 185 & 1 \\
3 & 1020 & 1040 & 370 & 0 \\
\hline
\end{tabular}
\end{center}
\end{examplebox}

\section{分层路由与自治系统}
\subsection{自治系统 (AS)}
自治系统是在单一的技术管理下的一组路由器，这些路由器使用一种自治系统内部的路由选择协议和共同的度量以确定分组在该自治系统内的路由。
\begin{itemize}[leftmargin=2em,itemsep=0pt]
    \item \textbf{作用}：屏蔽内部网络细节，简化路由管理，提高路由搜索效率。
    \item \textbf{AS 标识}：全球唯一的 16 位或 32 位 AS 编号。
\end{itemize}

\subsection{路由选择协议分类}
\begin{itemize}[leftmargin=2em,itemsep=0pt]
    \item \textbf{内部网关协议 (IGP)}：在一个自治系统内部使用的路由选择协议（如 RIP、OSPF）。
    \item \textbf{外部网关协议 (EGP)}：在不同自治系统之间使用的路由选择协议（如 BGP）。
\end{itemize}

\section{路由选择协议详解}
\subsection{RIP协议}
基于距离向量的内部网关协议。

\textbf{核心规则}：
\begin{itemize}[leftmargin=2em,itemsep=0pt]
    \item 以"跳数"作为度量标准
    \item 最大跳数为15，跳数=16表示不可达
    \item 每30秒向相邻路由器广播路由表
\end{itemize}

\subsection{OSPF协议}
基于链路状态的内部网关协议。

\textbf{核心特点}：
\begin{itemize}[leftmargin=2em,itemsep=0pt]
    \item 使用Dijkstra最短路径算法
    \item 以"链路代价"作为度量标准
    \item 无路由环路问题，收敛速度快
\end{itemize}

\subsection{BGP协议}
自治系统之间使用的外部网关协议。

\begin{notebox}
\textbf{路由协议的封装方式}：
\begin{itemize}[itemsep=0pt]
    \item RIP：封装在UDP（端口520）中
    \item OSPF：直接封装在IP（协议号89）中
    \item BGP：封装在TCP（端口179）中
\end{itemize}
\end{notebox}

\section{IP数据报转发}
\subsection{转发规则}
\begin{enumerate}[leftmargin=2em,itemsep=0pt]
    \item 计算自身网络地址：自身IP AND 子网掩码
    \item 计算目的网络地址：目的IP AND 自身子网掩码
    \item 若相同：直接交付；若不同：间接交付（发给网关）
\end{enumerate}

\begin{examplebox}[例题4.6 IP转发判断]
\textbf{题目}：主机A的IP为192.168.1.100/24，要发送数据给192.168.1.50和192.168.2.50。
\tcblower
A的网络地址：192.168.1.0

目的192.168.1.50的网络地址：192.168.1.0 → \textbf{直接交付}

目的192.168.2.50的网络地址：192.168.2.0 → \textbf{间接交付}
\end{examplebox}

\section{ICMP协议}
ICMP用于在主机和路由器之间传递控制信息和错误报告。

\subsection{差错报告报文}
\begin{itemize}[leftmargin=2em,itemsep=0pt]
    \item \textbf{终点不可达}：路由器无法转发数据报
    \item \textbf{源点抑制}：因拥塞丢弃数据报
    \item \textbf{时间超过}：TTL减为0时丢弃
    \item \textbf{参数问题}：IP首部字段有错误
    \item \textbf{改变路由}：通知源主机存在更优路径
\end{itemize}

\subsection{询问报文}
\begin{itemize}[leftmargin=2em,itemsep=0pt]
    \item \textbf{回送请求/回答}：ping命令使用
    \item \textbf{时间戳请求/回答}：用于时钟同步
\end{itemize}

\subsection{ICMP应用}
\textbf{ping}：使用ICMP回送请求和回答，测试连通性

\textbf{traceroute}：利用TTL递增和ICMP时间超过报文，逐跳发现路径

\section{VPN与NAT}
\subsection{虚拟专用网VPN}
利用公用因特网作为专用网之间的通信载体，使用IP隧道技术封装内部数据报。

\subsection{网络地址转换NAT}
允许多个私有IP地址共享一个公网IP地址。

\textbf{NAPT}：使用IP地址+端口号进行映射，一个公网IP可支持多台内网主机。

\begin{examplebox}[例题4.7 网络设备与广播域]
\textbf{题目}：下列网络设备中，能够抑制广播风暴的是？\\
Ⅰ 中继器 \quad Ⅱ 集线器 \quad Ⅲ 网桥 \quad Ⅳ 路由器
\tcblower
\begin{itemize}[itemsep=0pt]
    \item 中继器、集线器（物理层）：不隔离广播域
    \item 网桥（数据链路层）：不隔离广播域
    \item 路由器（网络层）：隔离广播域
\end{itemize}
\textbf{答案}：仅Ⅳ（路由器）
\end{examplebox}

\section{IPv6协议}
IPv6是下一代互联网协议，解决IPv4地址枯竭问题。

\subsection{IPv6地址特点}
\begin{itemize}[leftmargin=2em,itemsep=0pt]
    \item \textbf{地址长度}：128位（IPv4为32位）
    \item \textbf{地址数量}：$2^{128}$个，几乎无限
    \item \textbf{表示方法}：冒号十六进制，如\texttt{2001:0db8:85a3:0000:0000:8a2e:0370:7334}
    \item \textbf{简化规则}：前导0可省略，连续的0可用::表示（只能用一次）
\end{itemize}

\subsection{IPv6与IPv4的主要区别}
\begin{table}[htbp]
\centering\small
\begin{tabular}{lcc}
\toprule
\textbf{特性} & \textbf{IPv4} & \textbf{IPv6} \\
\midrule
地址长度 & 32位 & 128位 \\
首部长度 & 20-60字节 & 固定40字节 \\
首部校验和 & 有 & 无 \\
分片 & 路由器可分片 & 仅源主机分片 \\
广播 & 支持 & 不支持（用多播替代） \\
\bottomrule
\end{tabular}
\end{table}

\subsection{IPv6 over IPv4 (隧道技术)}
隧道技术允许 IPv6 分组通过现有的 IPv4 网络基础设施进行传输。
\begin{itemize}[leftmargin=2em,itemsep=0pt]
    \item \textbf{工作原理}：将整个 IPv6 数据报封装在 IPv4 数据报的数据部分中。
    \item \textbf{封装过程}：入口路由器将 IPv6 报文当作载荷，加上 IPv4 首部（协议号设为 41）。
    \item \textbf{解封过程}：出口路由器剥离 IPv4 首部，还原出原始的 IPv6 报文。
\end{itemize}

\section{路由表与 MAC 地址表的关系}
理解 IP 头部和数据链路 (DL) 头部地址在传输过程中的变化逻辑。
\begin{itemize}[leftmargin=2em,itemsep=0pt]
    \item \textbf{不变项}：在整个传输路径中，源 IP 和目的 IP 地址通常保持不变（除非经过 NAT）。
    \item \textbf{变化项}：源 MAC 和目的 MAC 地址在经过每一个路由器（跨越网段）时都会发生改变。
\end{itemize}

\begin{examplebox}[例题4.11 地址变化过程]
\textbf{场景}：主机A $\to$ 路由器R $\to$ 主机B。
\tcblower
\begin{enumerate}[itemsep=0pt]
    \item \textbf{A $\to$ R 段}：
    \begin{itemize}[itemsep=0pt]
        \item IP头部：(源A, 目的B) 
        \item DL头部：(源A-MAC, 目的R-MAC)
    \end{itemize}
    \item \textbf{R $\to$ B 段}：
    \begin{itemize}[itemsep=0pt]
        \item IP头部：(源A, 目的B) —— \textit{保持不变}
        \item DL头部：(源R-MAC, 目的B-MAC) —— \textit{发生重写}
    \end{itemize}
\end{enumerate}
\end{examplebox}

\begin{examplebox}[综合题4.11b IPv4首部分析]
\textbf{题目}：如下所示是捕获到的一个 IPv4 数据报首部，请分析各字段含义。
\begin{lstlisting}[basicstyle=\ttfamily\footnotesize]
Internet Protocol Version 4, Src: 192.168.1.100, Dst: 220.181.38.148
    0100 .... = Version: 4
    .... 0101 = Header Length: 20 bytes (5)
    Total Length: 52
    Identification: 0xab3c (43836)
    Flags: 0x40, Don't fragment
    Time to Live: 64
    Protocol: TCP (6)
    Header Checksum: 0x0000 [validation disabled]
    Source Address: 192.168.1.100
    Destination Address: 220.181.38.148
\end{lstlisting}
\tcblower
\textbf{字段分析}：
\begin{enumerate}[itemsep=0pt]
    \item \textbf{Version}：4，表示是 IPv4 协议。
    \item \textbf{Header Length}：20字节，这是最小值，表示没有选项字段。
    \item \textbf{Total Length}：52字节，表示首部+数据的总大小。
    \item \textbf{Identification}：0xab3c，分片时用于标识属于同一个数据报。
    \item \textbf{Flags}：Don't fragment，表示该报文不允许分片。
    \item \textbf{TTL}：64，表示报文在网络中最多还可以经过64个路由器。
    \item \textbf{Protocol}：TCP (6)，表示上层数据是使用 TCP 协议封装的。
    \item \textbf{Source/Destination}：显示发送方和接收方的 IP 地址。
\end{enumerate}
\end{examplebox}

\section{网络层综合例题}

\begin{examplebox}[例题4.8 子网划分综合]
\textbf{题目}：某单位分配到地址块\texttt{202.118.1.0/24}，需要划分为5个子网，子网1需要50台主机，子网2需要30台主机，子网3-5各需要10台主机。请设计地址分配方案。
\tcblower
采用VLSM变长子网掩码：

\textbf{子网1}（50台）：需要6位主机号（$2^6-2=62\geq50$）
\begin{itemize}[itemsep=0pt]
    \item 地址块：\texttt{202.118.1.0/26}
    \item 可用范围：\texttt{202.118.1.1 - 202.118.1.62}
\end{itemize}

\textbf{子网2}（30台）：需要5位主机号（$2^5-2=30\geq30$）
\begin{itemize}[itemsep=0pt]
    \item 地址块：\texttt{202.118.1.64/27}
    \item 可用范围：\texttt{202.118.1.65 - 202.118.1.94}
\end{itemize}

\textbf{子网3}（10台）：需要4位主机号（$2^4-2=14\geq10$）
\begin{itemize}[itemsep=0pt]
    \item 地址块：\texttt{202.118.1.96/28}
    \item 可用范围：\texttt{202.118.1.97 - 202.118.1.110}
\end{itemize}

\textbf{子网4}（10台）：
\begin{itemize}[itemsep=0pt]
    \item 地址块：\texttt{202.118.1.112/28}
    \item 可用范围：\texttt{202.118.1.113 - 202.118.1.126}
\end{itemize}

\textbf{子网5}（10台）：
\begin{itemize}[itemsep=0pt]
    \item 地址块：\texttt{202.118.1.128/28}
    \item 可用范围：\texttt{202.118.1.129 - 202.118.1.142}
\end{itemize}
\end{examplebox}

\begin{examplebox}[例题4.9 路由表查找]
\textbf{题目}：路由器R的路由表如下，分析目的地址为\texttt{172.32.57.14}的数据报如何转发。

\begin{center}\small
\begin{tabular}{|c|c|c|}
\hline
目的网络 & 子网掩码 & 下一跳 \\
\hline
172.32.0.0 & 255.255.0.0 & 接口0 \\
172.32.32.0 & 255.255.224.0 & 接口1 \\
172.32.56.0 & 255.255.252.0 & 接口2 \\
0.0.0.0 & 0.0.0.0 & 192.168.1.1 \\
\hline
\end{tabular}
\end{center}
\tcblower
对目的地址\texttt{172.32.57.14}逐条匹配：

\textbf{条目1}（/16）：172.32.57.14 AND 255.255.0.0 = 172.32.0.0 ✓

\textbf{条目2}（/19）：172.32.57.14 AND 255.255.224.0
\begin{itemize}[itemsep=0pt]
    \item 57 = 00111001，224 = 11100000，AND = 00100000 = 32
    \item 结果：172.32.32.0 ✓
\end{itemize}

\textbf{条目3}（/22）：172.32.57.14 AND 255.255.252.0
\begin{itemize}[itemsep=0pt]
    \item 57 = 00111001，252 = 11111100，AND = 00111000 = 56
    \item 结果：172.32.56.0 ✓
\end{itemize}

\textbf{条目4}（/0）：默认路由，总是匹配 ✓

根据\textbf{最长前缀匹配}原则，选择/22（条目3），转发到\textbf{接口2}
\end{examplebox}

\begin{examplebox}[例题4.10 IP路由器功能]
\textbf{题目}：下列关于IP路由器的描述，正确的是？\\
Ⅰ 运行路由协议，设置路由表\\
Ⅱ 监测到拥塞时，合理丢弃IP分组\\
Ⅲ 对收到的IP分组进行差错校验，确保传输的IP分组不丢失\\
Ⅳ 根据收到的IP分组的目的IP地址，将其转发到合适的输出线路上
\tcblower
\begin{itemize}[itemsep=0pt]
    \item Ⅰ正确：路由器运行路由协议维护路由表
    \item Ⅱ正确：拥塞时路由器会丢弃分组
    \item Ⅲ错误：IP层不负责可靠传输，仅校验首部，出错则丢弃
    \item Ⅳ正确：根据目的地址查表转发
\end{itemize}

\textbf{答案}：Ⅰ、Ⅱ、Ⅳ
\end{examplebox}

\section{本章习题}

\begin{examplebox}[填空题]
\textbf{1.} IPv4向IPv6进行过渡使用的两种策略是\underline{双协议栈}和\underline{隧道技术}。

\textbf{2.} 路由器数据层面的交换结构类型有\underline{经内存交换}、\underline{经总线交换}、\underline{经互联网络交换}。

\textbf{3.} B类IP地址网络允许容纳的最大主机数为\underline{65534}（$2^{16}-2$）。

\textbf{4.} 为了解决IP地址耗尽的问题，治本的措施是\underline{采用IPv6}。
\end{examplebox}

\begin{examplebox}[选择题]
\textbf{1.} 一个局域网与广域网互连，应该选用的互联设备为（\textbf{A}）\\
A. 路由器 \quad B. 网桥 \quad C. HUB \quad D. 交换机

\textbf{2.} 下列哪种协议不属于网络层协议（\textbf{D}）\\
A. ICMP \quad B. ARP \quad C. IP \quad D. DNS

\textbf{3.} 下面哪一种动态路由选择协议用作连接不同的互联网服务商自治系统间的协议（\textbf{B}）\\
A. OSPF \quad B. BGP \quad C. RIP \quad D. IS-IS

\textbf{4.} 为了解决IP地址耗尽的问题，可以采用以下的一些措施，其中治本的是（\textbf{A}）\\
A. 采用IPv6 \quad B. 划分子网 \quad C. 采用CIDR \quad D. 采用NAT

\textbf{5.} 在CRC的多项式表示方法中，若G(x)=101101，则多项式为（\textbf{C}）\\
A. $x^4+x^2+1$ \quad B. $x^5+x^4+x^2$ \quad C. $x^5+x^3+x^2+1$ \quad D. $x^5+x^3+1$
\end{examplebox}

\begin{examplebox}[计算题4.11 子网划分]
\textbf{题目}：假设某单位分配到一个地址块126.20.12.0/22。现要进一步划分为4个一样大的子网。试问：(1)每个子网的网络前缀有多长？(2)每个子网中有多少个地址？可以分配给主机使用的地址有多少个？(3)每个子网的地址块是什么？
\tcblower
(1) 原地址块前缀为/22，划分为4个子网需增加2位前缀，每个子网的网络前缀为\textbf{/24}

(2) 每个子网的地址数为$2^{32-24}=2^8=\textbf{256}$个，可分配给主机的地址数为$256-2=\textbf{254}$个（去除网络地址和广播地址）

(3) 每个子网的地址块：
\begin{itemize}[itemsep=0pt]
    \item 子网1：126.20.12.0/24
    \item 子网2：126.20.13.0/24
    \item 子网3：126.20.14.0/24
    \item 子网4：126.20.15.0/24
\end{itemize}
\end{examplebox}

\begin{examplebox}[计算题4.12 路由表转发]
\textbf{题目}：设某路由器建立了如下路由表，现共收到4个分组，其目的地址分别为：(1)128.96.39.10 (2)128.96.40.12 (3)128.96.40.151 (4)192.4.153.17，求各分组的转发路径。

\begin{center}\small
\begin{tabular}{|c|c|c|}
\hline
目的网络 & 子网掩码 & 下一跳 \\
\hline
128.96.39.0 & 255.255.255.128 & 接口m0 \\
128.96.39.128 & 255.255.255.128 & 接口m1 \\
128.96.40.0 & 255.255.255.128 & R2 \\
192.4.153.0 & 255.255.255.192 & R3 \\
0.0.0.0（默认） & 0.0.0.0 & R4 \\
\hline
\end{tabular}
\end{center}
\tcblower
(1) \textbf{128.96.39.10}：与128.96.39.0/25匹配，下一跳为\textbf{接口m0}

(2) \textbf{128.96.40.12}：与128.96.40.0/25匹配，下一跳为\textbf{R2}

(3) \textbf{128.96.40.151}：151>127，不匹配128.96.40.0/25，走默认路由，下一跳为\textbf{R4}

(4) \textbf{192.4.153.17}：与192.4.153.0/26匹配，下一跳为\textbf{R3}
\end{examplebox}

\begin{examplebox}[计算题4.13 IP数据报分片]
\textbf{题目}：原始IP数据报总长度为6578字节，标识为942，假设该IP分组需要通过以太网帧传输（MTU=1500字节），说明如何进行IP数据包的分片？每个分片的总长度是多少？每片的MF标志是多少？片偏移是多少？
\tcblower
以太网MTU=1500字节，IP头默认20字节，每个分片的数据部分最大为1480字节

总数据长度 = 6578 - 20 = 6558字节

需分片数 = $\lceil 6558/1480 \rceil$ = 5片

\textbf{分片详情}：
\begin{center}\small
\begin{tabular}{|c|c|c|c|c|}
\hline
分片 & 数据长度 & 总长度 & MF & 片偏移 \\
\hline
1 & 1480 & 1500 & 1 & 0 \\
2 & 1480 & 1500 & 1 & 185 (1480/8) \\
3 & 1480 & 1500 & 1 & 370 (2960/8) \\
4 & 1480 & 1500 & 1 & 555 (4440/8) \\
5 & 638 & 658 & 0 & 740 (5920/8) \\
\hline
\end{tabular}
\end{center}
\end{examplebox}

\begin{examplebox}[综合题：ARP协议分析]
\textbf{题目}：分析下图所示数据包各个字段的含义是什么？并描述ARP协议的工作原理。

\begin{lstlisting}[basicstyle=\ttfamily\footnotesize]
Address Resolution Protocol (reply)
  Hardware type: Ethernet (1)
  Protocol type: IPv4 (0x0800)
  Hardware size: 6
  Protocol size: 4
  Opcode: reply (2)
  Sender MAC address: 00:0f:e2:e1:c9:06
  Sender IP address: 172.28.31.254
  Target MAC address: 74:27:ea:4f:47:d2
  Target IP address: 172.20.28.40
\end{lstlisting}
\tcblower
\textbf{字段含义}：
\begin{enumerate}[itemsep=0pt]
    \item Hardware type: 硬件类型，1表示以太网
    \item Protocol type: 协议类型，0x0800表示IPv4
    \item Hardware size: 硬件地址长度，以太网MAC地址为6字节
    \item Protocol size: 协议地址长度，IPv4地址为4字节
    \item Opcode: 操作码，1为请求，2为应答
    \item Sender MAC/IP: 发送方的MAC地址和IP地址
    \item Target MAC/IP: 目标方的MAC地址和IP地址
\end{enumerate}

\textbf{ARP工作原理}：
\begin{enumerate}[itemsep=0pt]
    \item 主机检查本地ARP缓存
    \item 若没有，广播ARP请求（目标MAC为FF-FF-FF-FF-FF-FF）
    \item 目标主机收到请求后，单播回复ARP应答
    \item 发送方将映射存入ARP缓存
\end{enumerate}
\end{examplebox}

\begin{examplebox}[选择题（补充）]
\textbf{6.} 下面可以分配给计算机的有效IP地址是（\textbf{A}）\\
A. 202.221.130.1 \quad B. 127.192.33.41 \quad C. 191.0.130.0 \quad D. 283.192.33.1

\textit{解析}：B是环回地址，C主机号全0是网络地址，D超出范围。

\textbf{7.} 以下各项中，不是数据报操作特点的是（\textbf{C}）\\
A. 每个分组自身携带有足够的信息，它的传送是被单独处理的\\
B. 在整个传送过程中，不需建立虚电路\\
C. 使所有分组按顺序到达目的端系统\\
D. 网络节点要为每个分组做出路由选择

\textbf{8.} 下列选项中，不属于物理层接口规范定义范畴的是（\textbf{A}）\\
A. 物理地址 \quad B. 接口形状 \quad C. 引脚功能 \quad D. 信号电平

\textbf{9.} 双绞线传输介质是把两根导线绞在一起，这样可以减少（\textbf{D}）\\
A. 外界信号的干扰 \quad B. 信号传输时的衰减\\
C. 信号向外泄漏 \quad D. 信号之间的相互串扰
\end{examplebox}

\begin{examplebox}[计算题4.14 子网划分（变形）]
\textbf{题目}：某公司申请到地址块202.113.79.0/24，需要划分为若干子网：研发部60台主机，市场部28台主机，财务部12台主机，行政部6台主机。设计地址分配方案。
\tcblower
采用VLSM，按需求从大到小分配：

\textbf{研发部}（60台）：需要6位主机号（$2^6-2=62\geq60$）
\begin{itemize}[itemsep=0pt]
    \item 地址块：202.113.79.0/26
    \item 可用范围：202.113.79.1 - 202.113.79.62
\end{itemize}

\textbf{市场部}（28台）：需要5位主机号（$2^5-2=30\geq28$）
\begin{itemize}[itemsep=0pt]
    \item 地址块：202.113.79.64/27
    \item 可用范围：202.113.79.65 - 202.113.79.94
\end{itemize}

\textbf{财务部}（12台）：需要4位主机号（$2^4-2=14\geq12$）
\begin{itemize}[itemsep=0pt]
    \item 地址块：202.113.79.96/28
    \item 可用范围：202.113.79.97 - 202.113.79.110
\end{itemize}

\textbf{行政部}（6台）：需要4位主机号（$2^4-2=14\geq6$）
\begin{itemize}[itemsep=0pt]
    \item 地址块：202.113.79.112/28
    \item 可用范围：202.113.79.113 - 202.113.79.126
\end{itemize}
\end{examplebox}

\begin{examplebox}[计算题4.15 路由聚合]
\textbf{题目}：将以下4个网络聚合为一个超网：192.168.16.0/24、192.168.17.0/24、192.168.18.0/24、192.168.19.0/24。
\tcblower
第三字节二进制：
\begin{itemize}[itemsep=0pt]
    \item 16 = 00010000
    \item 17 = 00010001
    \item 18 = 00010010
    \item 19 = 00010011
\end{itemize}

共同前缀：000100（前6位相同）

聚合后前缀长度 = 16 + 6 = 22

\textbf{聚合结果}：192.168.16.0/22
\end{examplebox}

\begin{examplebox}[计算题4.16 IP分片（变形）]
\textbf{题目}：一个IP数据报总长度为3820字节（首部20字节），需要通过MTU=800字节的链路传输。求分片数量和每个分片的详细信息。
\tcblower
数据部分 = 3820 - 20 = 3800字节

每个分片最大数据 = 800 - 20 = 780字节

但780不是8的倍数，需要调整为776字节（$776 = 97 \times 8$）

分片数 = $\lceil 3800/776 \rceil$ = 5片

\textbf{分片详情}：
\begin{center}\small
\begin{tabular}{|c|c|c|c|c|}
\hline
分片 & 数据长度 & 总长度 & MF & 片偏移 \\
\hline
1 & 776 & 796 & 1 & 0 \\
2 & 776 & 796 & 1 & 97 \\
3 & 776 & 796 & 1 & 194 \\
4 & 776 & 796 & 1 & 291 \\
5 & 696 & 716 & 0 & 388 \\
\hline
\end{tabular}
\end{center}
\end{examplebox}

\begin{examplebox}[简答题：RIP与OSPF对比]
\textbf{题目}：比较RIP和OSPF路由协议的主要区别。
\tcblower
\begin{center}\small
\begin{tabular}{|l|c|c|}
\hline
\textbf{特性} & \textbf{RIP} & \textbf{OSPF} \\
\hline
算法类型 & 距离向量 & 链路状态 \\
度量标准 & 跳数 & 链路代价 \\
最大跳数 & 15 & 无限制 \\
收敛速度 & 慢 & 快 \\
路由环路 & 可能产生 & 不会产生 \\
更新方式 & 周期性广播 & 触发式多播 \\
封装协议 & UDP & IP \\
适用规模 & 小型网络 & 大型网络 \\
\hline
\end{tabular}
\end{center}
\end{examplebox}
